\section{An unconditional Gram improvement and the converse trace implication}
\label{sec:trace-unconditional-converse}

The forward rational trace theorem above assumes \(\mathsf A\).
Its converse must not reuse the inverse-Gram estimate AC-T28,
which also assumes \(\mathsf A\).
The following weaker, unconditional improvement is enough for
the converse.  It uses cancellation in a finite Mobius inverse,
and does not identify a trace limit without a trace hypothesis.

\begin{lemma}[AC-T49: the finite Mobius inverse at polynomial frequencies]
\label{AC-T49}
Fix a Dirichlet character \(\chi\), including the character
of conductor one.  For every fixed \(K,B>0\),
\[
 b_{\chi,N}(\tau)
   :=\sum_{n\leq N}\frac{\mu(n)\chi(n)}{n^{1/2+i\tau}}
 \quad\hbox{satisfies}\quad
 \sup_{|\tau|\leq N^K}|b_{\chi,N}(\tau)|
       \leq C_{\chi,K,B}\frac{\sqrt N}{(\log N)^B}
 \tag{UC.1}
\]
for all sufficiently large \(N\).
On the whole real line one also has
\[
                         |b_{\chi,N}(\tau)|\leq2\sqrt N .
 \tag{UC.2}
\]
Both estimates are unconditional.
\end{lemma}
\begin{proof}
We use the published classical Dirichlet polynomial estimate
in Matom\"aki--Shao--Tao--Ter\"av\"ainen,
\emph{Higher uniformity of arithmetic functions in short
intervals I. All intervals}, \emph{Forum Math. Pi} 11
(2023), e29, Section 3.4, Lemma 3.9(ii), printed p.~46;
its proof is on pp.~46--47,
\url{https://doi.org/10.1017/fmp.2023.28}.
For fixed \(0<\alpha\leq1\), that lemma bounds
\[
 \sup_{|t|\leq X}\sup_{I\subset[X^\alpha,2X^\alpha]}
 \left|\sum_{n\in I}\frac{\mu(n)\chi(n)}{n^{1/2+it}}\right|
       \ll_{\alpha,A}\frac{X^{\alpha/2}}{(\log X)^A}
 \tag{UC.3}
\]
when the modulus is at most \((\log X)^A\).
This is its case \(r=1\).
The source obtains this part from the classical zero-free-region
contour estimate, without an RH hypothesis.  That arithmetic
theorem is an explicitly cited background input here.

Increase \(K\) to \(\max(K,1)\).
Discard first the integers \(n\leq2N^{1/2}\), and partition
the remaining integers into dyadic intervals
\(I\subset[Y,2Y]\), with \(Y\geq2N^{1/2}\), allowing
a final truncated interval.
For each such interval set
\[
                       X=Y^{2K},\qquad \alpha=(2K)^{-1}.
\]
Then \(X^\alpha=Y\), \(X\geq N^K\), and the fixed modulus
is at most \((\log X)^B\) for all sufficiently large \(N\).
Equation (UC.3), with \(A=B\), bounds each dyadic piece by
\(C_{\chi,K,B}\sqrt Y/(\log N)^B\).
The sum of their square roots is at most a constant times
\(\sqrt N\), since their lengths form a geometric sequence.
The initial segment contributes at most
\(2\sqrt{2N^{1/2}}=O(N^{1/4})\).
This is absorbed by \(\sqrt N/(\log N)^B\), proving (UC.1).
Finally \(|\mu(n)\chi(n)|\leq1\) and
\(\sum_{n\leq N}n^{-1/2}\leq2\sqrt N\) prove (UC.2).
\end{proof}

\begin{theorem}[AC-T50: unconditional Gram control sufficient for the tested defect]
\label{AC-T50}
Fix a primitive character \(\chi\), and use either the
original full parity-index prolate family of AC-T29 or
one of the actual fixed-\(S\), fixed-\(\eta\) angle families
of AC-T48.  In the principal sector take its exact
two-constraint raw subspace
\(\ker I\cap\ker\operatorname{ev}_0\).
Use ordinary positive-half-line raw orthonormal coordinates
on the chosen finite subspace.
Let \(E_\Lambda\) be its global theta column operator,
corrected as in AC-T42 in the principal sector, and put
\[
       G_\Lambda=E_\Lambda^*E_\Lambda,\qquad
       Z_\Lambda=E_\Lambda G_\Lambda^{-1/2},\qquad
       Q_\Lambda=Z_\Lambda Z_\Lambda^*,\qquad
       c=\Lambda^2/q .
\]
For every fixed \(B>0\), unconditionally,
\[
 \|G_\Lambda^{-1}\|
       \leq C_{\chi,S,\eta,B}\frac{c}{(\log c)^{2B}}
 \tag{UC.4}
\]
for all sufficiently large \(\Lambda\).
The \(S,\eta\) dependence is omitted for the original family.

Let \(S_\Lambda=1_{[\Lambda^{-1},\Lambda]}\) on the
global target, and let
\(\mathscr G=C_c^\infty(\mathbb R_+^\times)+\mathbb C k\),
where \(k(e^t)=\frac12e^{-|t|}\).
For every \(g\in\mathscr G\) and every \(B>0\),
\[
 \|V(g)(1-S_\Lambda)Z_\Lambda\|_2^2
        \leq C_{\chi,S,\eta,g,B}(\log\Lambda)^{-2B}.
 \tag{UC.5}
\]
Here the norm is Hilbert--Schmidt.
Consequently, for every \(h\in C_c^\infty(\mathbb R_+^\times)\)
and every \(M>0\),
\[
 \left|\operatorname{Tr}
       \bigl((Q_\Lambda-S_\Lambda Q_\Lambda S_\Lambda)V(h)\bigr)
       \right|
       \leq C_{\chi,S,\eta,h,M}(\log\Lambda)^{-M}.
 \tag{UC.6}
\]
In particular, for \(g\in C_c^\infty(\mathbb R_+^\times)\),
\[
 \liminf_{\Lambda\to\infty}
 \left\{2\log\Lambda\,\|g\|_2^2
        -\operatorname{Tr}(Q_\Lambda V(g*g^*))\right\}\geq0 .
 \tag{UC.7}
\]
None of (UC.4)--(UC.7) assumes \(\mathsf A\) or the
existence of a renormalized trace limit.
\end{theorem}
\begin{proof}
\emph{1. Unconditional raw estimates.}
Write \(R=\Lambda/q\), and let \(\mathcal R_\Lambda\)
denote the selected raw subspace, with the two constraints
in the principal case.
There are fixed constants \(d,C>0\) such that, for all
\(\psi\in\mathcal R_\Lambda\) and all sufficiently large
\(\Lambda\),
\[
 \begin{split}
 &\operatorname{supp}\psi\subset[0,\Lambda],\qquad
 a=d/R,\qquad
 \int_a^\Lambda|\psi(u)|^2du\geq\tfrac12\|\psi\|_2^2,\\
 &\|\psi\|_{H_D^{3/8}}\leq Cc^{3/8}\|\psi\|_2 .
 \end{split}
 \tag{UC.8}
\]
For the original family, (TB.27)--(TB.28) prove the mass
bound, and (TB.33)--(TB.34), applied to
\(e^{t/2}\psi(e^t)\), prove the Sobolev bound.
The original concentration values are uniformly bounded
below by (TB.49); this scalar fact is unconditional.
For the actual angle family these are (AD.20) and (AD.14),
whose proofs use only its band extension, the fixed
transported metric and the fractional cutoff estimate.
In particular these portions of AC-T48 do not use its
later RH-dependent inverse-Gram step.
Both estimates pass to a constrained subspace.

Let \(f(t)=e^{t/2}\psi(e^t)\), and write
\(e(t)=E_\chi\psi(e^t)\).
The unconditional critical-line multiplier bound (TB.16),
or its principal version (TP.3), gives
\[
              \|e\|_{H^{1/16}}
                   \leq C_\chi\|\psi\|_{H_D^{3/8}}
                   \leq Cc^{3/8}\|\psi\|_2 .
 \tag{UC.9}
\]
The exponents add exactly:
\(1/16+5/16=3/8<1/2\).
The principal mean is zero, so in that case as well
\(e(t)=0\) for \(t>\log\Lambda\).

\emph{2. Apply the finite inverse before taking its norm.}
Let \(N=\lceil\Lambda/a\rceil\), so \(N\asymp c\), and set
\[
 (\mathcal B_Ne)(t)
       =\sum_{n\leq N}\frac{\mu(n)\chi(n)}{\sqrt n}
                                      e(t+\log n).
 \tag{UC.10}
\]
Finite Dirichlet convolution gives the exact identity
\[
              f(t)=\mathcal B_Ne(t),
            \qquad \log a\leq t\leq\log\Lambda .
 \tag{UC.11}
\]
Indeed the upper support makes all terms with
\(n>\Lambda e^{-t}\) zero.  Expanding the remaining theta
sums gives coefficient
\(\chi(k)\sum_{n\mid k}\mu(n)\) at \(\sqrt{e^t}\psi(ke^t)\).
It vanishes except at \(k=1\).
For the principal character the corrected theta sum is the
uncorrected sum on this subspace because \(I\psi=0\).
For the possibly nonsmooth angle vectors the same equality
holds in local \(L^2\): on every compact target interval
bounded away from zero there are only finitely many sampling
maps, each bounded from the compact raw \(L^2\) space.
Approximation in the fractional graph domain proves their
agreement with the closed theta operator.
In the principal case preserve the continuous raw mean
condition during this approximation.
Thus no pointwise regularity at an internal boundary is required.

In the plus Fourier convention, the multiplier of
\(\mathcal B_N\) is precisely \(b_{\chi,N}(\tau)\).
Split its action on \(e\) at \(|\tau|=N^{16}\).
Equations (UC.1)--(UC.2) and Plancherel give
\[
 \begin{split}
 \|\mathcal B_Ne\|_2
 &\leq C_{\chi,B}\frac{\sqrt N}{(\log N)^B}\|e\|_2
     +2\sqrt N\,N^{-1}\|e\|_{H^{1/16}}\\
 &\leq C_{\chi,S,\eta,B}
              \frac{\sqrt c}{(\log c)^B}\|e\|_2
        +C_{\chi,S,\eta}c^{-1/8}\|\psi\|_2 .
 \end{split}
 \tag{UC.12}
\]
The factor \(N^{-1}\) is
\((N^{16})^{-1/16}\), and the final power is
\(1/2-1+3/8=-1/8\).
By (UC.8) and (UC.11), the left side is at least
\(2^{-1/2}\|\psi\|_2\).
For sufficiently large \(c\), absorb the last term in
(UC.12).  Squaring the resulting estimate proves
\[
       \|E_\chi\psi\|_2^2
          \geq C_{\chi,S,\eta,B}^{-1}
                      \frac{(\log c)^{2B}}c\,\|\psi\|_2^2 .
 \tag{UC.13}
\]
In raw orthonormal coordinates this is exactly (UC.4).
No reciprocal bound for \(L\) on a shifted critical line
has occurred in this argument.

\emph{3. The actual smoothed tail and the trace.}
The unconditional parts of AC-T29, (TP.27), and
(AD.21)--(AD.23) give, for the respective families,
\[
            \|V(g)(1-S_\Lambda)E_\Lambda\|_2^2
                     \leq C_{\chi,S,\eta,g}\Lambda^{-2}.
 \tag{UC.14}
\]
For clarity, the principal argument here uses the exact
two constraints.  The upper theta tail is zero by
\(I\psi=0\).  Its band extension agrees at zero with
\(\psi\), so its lower polar tail is zero by
\(\psi(0)=0\).  Euler sum-minus-integral gives the
outside-time kernel estimate (TP.26).
For the fixed-\(S\) angle family, the finite divisor
expansion of \(1-A=Z^{-1}(1-P_\Lambda)Z\) gives (AD.23).
Those kernel and band-extension bounds are unconditional;
(AD.19) is not being used.

Compose (UC.14) with \(G_\Lambda^{-1/2}\) and use
\(\Lambda^{-2}c=1/q\).
This proves (UC.5).
Put \(Z_0=S_\Lambda Z_\Lambda\), \(Z_1=(1-S_\Lambda)Z_\Lambda\),
\(C_g=V(g)^*Z_1\), and \(D_g=V(g)^*Z_0\).
The exact identity (TB.54) and its energy estimate give
\[
 \begin{split}
 &\left|\operatorname{Tr}
  \bigl((Q_\Lambda-S_\Lambda Q_\Lambda S_\Lambda)
                                  V(g*g^*)\bigr)\right|\\
 &\hspace{12mm}\leq
 \|C_g\|_2^2+
       2\sqrt{2\log\Lambda}\,\|g\|_2\|C_g\|_2
       \leq C_{g,B}
           \bigl((\log\Lambda)^{-2B}
                    +(\log\Lambda)^{1/2-B}\bigr).
 \end{split}
 \tag{UC.15}
\]
The involution preserves \(\mathscr G\), so (UC.5) applies
to \(g^*\) as needed.
Since \(B\) is arbitrary, this is smaller than any prescribed
negative power of \(\log\Lambda\).
Polarization applies on \(\mathscr G\).
Every compact smooth logarithmic test \(h\) can be written
\(h=k*j^*\), where, in logarithmic coordinates,
\(j^*=(1-\partial_t^2)h\) is compact and smooth.
Indeed \(k\) is the fundamental solution of
\(1-\partial_t^2\).
Thus polarization of (UC.15) proves (UC.6) for every
original test.

Finally \(0\leq S_\Lambda Q_\Lambda S_\Lambda\leq S_\Lambda\).
The exact interval energy is
\[
         \operatorname{Tr}(V(g)^*S_\Lambda V(g))
                       =2\log\Lambda\,\|g\|_2^2 .
 \tag{UC.16}
\]
Subtracting the compressed positive contraction gives
a nonnegative trace.  Equation (UC.6) then proves (UC.7).
\end{proof}

\begin{theorem}[AC-T51: the rational canonical trace formula is equivalent to \(\mathsf A\)]
\label{AC-T51}
Use the canonical rational finite families, Hilbert metrics,
and centered representation \(\mathbb R\) of
AC-T45--AC-T47.  Let \(\mathcal W_{\rm geom}\) be the
centered Weil distribution of C99, Appendix II, Theorem 6,
equivalently CCM07, Theorem 6.1, in the normalization (T.1).
This geometric distribution is defined without \(\mathsf A\).
Put
\[
       \mathcal P(h)=\int_0^\infty h_1(r)
                           (r^{-1/2}+r^{1/2})\,\frac{dr}{r}.
 \tag{UC.17}
\]
Then the following two assertions are equivalent:
\[
 \begin{split}
 &\mathsf A;\\
 &\text{for every }h\in\mathcal D(C_{\mathbb Q}),\quad
 \lim_{\Lambda\to\infty}
 \left\{2\log\Lambda\,h(1,1)
       -\operatorname{Tr}(\mathbb Q_\Lambda\mathbb R(h))
       +\mathcal P(h)\right\}
                         =\mathcal W_{\rm geom}(h).
 \end{split}
 \tag{UC.18}
\]
The same equivalence holds after replacing, in each fixed
character sector, the original full-index family by the
fixed-\(S\), fixed-\(\eta\) family of AC-T48, with its
corresponding full principal completion.
Every \(S,\eta\) is held fixed as \(\Lambda\) tends to infinity.
\end{theorem}
\begin{proof}
The implication from \(\mathsf A\) is AC-T41,
AC-T44--AC-T46 and the exact source dictionary AC-T47;
AC-T48 gives the stated angle-family variant.
Their identification of the limiting zero distribution
uses \(\mathsf A\), as explicitly recorded in those proofs.

For the converse, assume the trace assertion in (UC.18)
for every original test, without assuming \(\mathsf A\).
The rank-two comparison (PC.4)--(PC.5) is unconditional.
It removes precisely the polar term \(\mathcal P(h)\).
For the fixed-\(S\) variant its proof uses the unconditional
fixed-profile and two-constraint approximation in step 2
of AC-T48, and is unchanged.
Thus the assumed limit equals the renormalized trace on
the direct sum of the nonprincipal theta targets and the
two-constraint principal theta target.

Take \(h=f*f^*\), with \(f\in\mathcal D(C_{\mathbb Q})\).
The compact-unit Fourier expansion has only finitely many
characters, and its coefficients satisfy
\[
                     h_\chi=f_\chi*f_\chi^* .
 \tag{UC.19}
\]
Apply (UC.7) to each of these finitely many sectors and
add the inequalities.  Fourier inversion gives
\(\sum_\chi\|f_\chi\|_2^2=(f*f^*)(1,1)\).
The assumed trace limit and the unconditional polar
comparison therefore imply
\[
                          \mathcal W_{\rm geom}(f*f^*)\geq0
               \qquad(f\in\mathcal D(C_{\mathbb Q})).
 \tag{UC.20}
\]
One may pass to the strong Schwartz domain in the cited
criterion without changing the test hypothesis.
Indeed CCM07 defines that domain in its equation
labelled \texttt{SCK} as
\(\bigcap_{\beta\in\mathbb R}|\cdot|^\beta\mathcal S(C_{\mathbb Q})\).
On each finite compact-unit quotient, take a smooth compact
logarithmic cutoff equal to one near zero and apply it at
scale \(j\).
The resulting \(f_j\in\mathcal D\) converge to \(f\) in
every exponentially weighted Schwartz seminorm: after
Leibniz differentiation, each discarded term lies in
\(\{|t|\geq c j\}\), where an extra exponential weight
makes its seminorm tend to zero.
Convolution and involution are continuous in these seminorms.
The explicit Weil formula is continuous there as well:
its prime terms are bounded by a seminorm with exponential
weight greater than \(1/2\), and its archimedean principal
value is bounded by finitely many weighted derivative
seminorms.
Thus (UC.20) holds on that strong Schwartz domain.

The converse Weil criterion in CCM07,
Proposition 6.2 (label \texttt{RHpositTr}) and
Corollary 6.3 (label \texttt{RHpositTrsharp}), with the
normalization checked in AC-T02, now gives RH for every
rational unitary Gr\"ossencharacter.
By AC-T01 this is exactly \(\mathsf A\).
All estimates used to pass from the assumed trace to
(UC.20) are unconditional; in particular AC-T28,
the half-line absorption under \(\mathsf A\), and the
upper zero-measure theorem have not been used in this
direction.
\end{proof}

\begin{remark}[What the unconditional improvement does and does not prove]
AC-T50 proves asymptotic positivity of the actual finite
cutoff defect and its tested comparison with a positive
window contraction.  It does not prove, unconditionally,
that this defect has the Weil distribution as its limit.
That identification is the hypothesis in the converse
half of AC-T51, and is the conclusion of the forward
half under \(\mathsf A\).
Consequently AC-T51 proves an equivalence, not an
unconditional proof of \(\mathsf A\).
The field in this theorem is \(\mathbb Q\), a number field.
It includes all of its unitary characters, not merely the
zeta sector.  No extension to all characters of an
arbitrary number field, no growing-conductor uniformity,
and no growing-\(S\) cutoff limit is asserted here.
\end{remark}
